\section{Conclusion}
\label{sec:conclusion}

In this work, we introduced MAYA (Manifold Alignment Yielding Accuracy), a gradient-free framework for online class-incremental learning under the strict constraints of frozen, pre-trained visual backbones. We identified that the primary driver of catastrophic failure in zero-buffer CIL is the severe geometric bias inherent to pre-trained feature spaces. To resolve this, MAYA employs a dual-stage approach: an analytic projection to an Equiangular Tight Frame (ETF) to globally unwarp and separate class manifolds, coupled with a streaming Episodic Graph to map the precise, multi-modal topography within those aligned spaces. 

Empirically, we demonstrated that MAYA bridges the critical performance gap between naive prototype classifiers and computationally expensive gradient-based replay methods. By explicitly avoiding the assumption of homoscedasticity, MAYA achieved state-of-the-art performance on topologically complex, out-of-distribution datasets like ObjectNet (+28.2\% over NCM), matching or exceeding the strict statistical bounds of highly specialized analytic methods like Deep SLDA and Exact ACL.

Crucially, the architecture of MAYA draws heavy inspiration from biological memory consolidation. The human brain does not incorporate new information by computing exact, global rank-1 matrix updates to a monolithic, dense neural network (as in statistical covariance methods). Instead, the mammalian hippocampus relies on the rapid formation of localized, sparse episodic memory traces—mirroring MAYA's dynamic node construction. Over time, these local traces interact with a broader, structurally stable semantic space (analogous to our global ETF projection). By utilizing an explicit Episodic Graph rather than a "black-box" statistical accumulator, MAYA not only achieves high accuracy but retains the biological advantages of interpretability and concept localization, allowing for the discrete unlearning or manipulation of specific memories without inducing catastrophic interference. 

Future work will explore dynamically expanding the global ETF dimension to accommodate theoretically infinite class streams, and extending the Episodic Graph to model smooth temporal shifts in streaming video data.
